\documentclass{exam}
\usepackage{../../mypackages}
\usepackage{../../macros}

\title{Interrogation N°2}
\author{N. Bancel}
\date{16 Octobre 2025}

\begin{document}

\textbf{Collège Lycée Suger}
\hfill
\textbf{Mathématiques} \\

\textbf{Année 2025-2026}
\hfill
\textbf{Classe de 6ème} \par

{\let\newpage\relax\maketitle}

\begin{center}
\textbf{\textcolor{red}{Durée : 1 heure. La calculatrice n'est pas autorisée}} \\
\textbf{\textcolor{red}{Une réponse donnée sans justification sera considérée comme fausse}} \\
Cette interrogation contient \numquestions\ questions, sur \numpages\ pages et est notée sur 20. \\
Une copie mal tenue sera pénalisée.

\end{center}

\section*{Notions et rappels du Chapitre 1}

\subsection*{Exercice 1 — Lecture et écriture des nombres (1.5 points)}

\begin{questions}
  \question[0.5] Écrire en lettres le nombre suivant : \quad $80\,431$.
  \question[1] Déterminer pour le nombre $2\,389,461$ :
  \begin{parts}
    \part Le nombre de dizaines.
    \part La partie décimale.
    \part Le chiffre des centaines.
    \part La partie décimale
  \end{parts}
\end{questions}

\subsection*{Exercice 2 — Fractions décimales et écritures décimales (2 points)}

\begin{questions}
  \question[0.5] Ecrire sous la forme d'une seule et même fraction décimale $45,981$
  \question[0.5] Ecrire sous la forme d'une décomposition par parties (entier + fraction décimale) $2,097$
  \question[0.5] Ecrire sous la forme d'une décomposition par rang ($... + \frac{...}{10} + \frac{...}{100} + \frac{...}{1000} + ...$) le nombre $34,0052$
  \question[0.5] Ecrire sous forme décimale $\frac{341}{1000}$
\end{questions}

\subsection*{Exercice 3 — Lecture sur une droite graduée (2 points)}

\begin{questions}
  \question[1] Donner l'abscisse de $T$ en écriture décimale. Justifier. La justification compte autant que la réponse.
\end{questions}

  \begin{center}
    \includegraphics[width=0.8\linewidth]{interro_2_01.jpg}
  \end{center}

\begin{questions}
  \question[1] Donner l'abscisse du point $W$ en écriture décimale. Justifier.
\end{questions}

    \begin{center}
    \includegraphics[width=0.8\linewidth]{interro_2_02.jpg}
  \end{center}

\subsection*{Exercice 4 — Comparaisons et rangement (2 points)}

\begin{questions}
  \question[0.5] Effectuer une décomposition du nombre $903\,821$
  \question[1.5] Classer les nombres suivants dans l'ordre \textbf{croissant} :  
  $87,09 \quad 82,2 \quad 82,92 \quad 87,764 \quad 82,467 \quad 82,764$.
\end{questions}

\section*{Applications du Chapitre 2}

\subsection*{Exercice 5 — Règles de priorité et calcul astucieux (5 points)}

\begin{questions}
  \question[2.5] Effectuer les calculs suivants en détaillant les étapes.
  \begin{parts}
    \part $A = 45 - (16 + 8) + 4$
    \part $B = 5 \times ((4 + 5) \times 2)$
    \part $C = (1 + 9) \times (14 - 3)$
    \part $D = 3 \times [10 - (2 + 3)] + 5$
  \end{parts}
\end{questions}

\begin{questions}
  \question[2.5] Effectuer les calculs suivants en faisant des regroupements astucieux - Mettre en lumière ces groupements astucieux (je voudrais voir comment vous avez regroupé les nombres) :
  \begin{parts}
    \part $D = 25 + 18 + 75 + 82$
    \part $E = 4,5 \times 2 \times 5$ 
    \part $F = 2,5 \times 5 \times 4 \times 177 \times 2$
    \part $C = 25 \times 4,7 \times 0,5 \times 4 \times 2$
  \end{parts}
\end{questions}

\subsection*{Exercice 6 - Calculs posés. (3.5 points)}

\textbf{En posant le calcul, déterminer la valeur de :}
\begin{questions}
  \question[1] $G = 7496,8 - 29,54$
  \question[1.5] $H = 49,3 \times 2,05$
  \question[1] $I = 496,891 + 25,72$
\end{questions}

\subsection*{Exercice 7 — Ordres de grandeur (1.5 point)}

\begin{questions}
  \question Donner un ordre de grandeur (et préciser à quel ordre de grandeur vous arrondissez) de :
  \begin{parts}
    \part[0.5] $J = 78,4 + 121,6$
    \part[0.5] $K = 2,98 \times 5,1$
    \part[0.5] $L = 1204,32 \times 98,8$
  \end{parts}

  
\end{questions}

\subsection*{Exercice 9 — Problème (2.5 point)}

\begin{questions}
  \question Je paye 100 feuilles à $0,05$ euros chacune, 1 classeur à 6 euros, et une règle à 2.4 euros. Je donne un billet de 50 euros.  
  Combien doit-on me rendre ?
\end{questions}

\end{document}
